\begin{definition}[5.9.9]
\label{IV.2.s5.9.9-fr}
Bajo las hipótesis de \textup{(5.9.8)}, se dice que $\mathcal F$ es $Z$-puro
(\textup{resp.} $Z$-cerrado) si el homomorfismo $\rho_{X/Z}$ es inyectivo
(\textup{resp.} biyectivo).
\end{definition}

Si $X=\operatorname{Spec}(A)$ es afín y $\mathcal F=\widetilde M$, donde $M$
es un $A$-módulo, se dirá que $M$ es $Z$-puro (resp. $Z$-cerrado) cuando
$\mathcal F$ sea $Z$-puro (resp. $Z$-cerrado).

Se dice que $\mathcal F$ es $Z$-\emph{puro} (resp. $Z$-\emph{cerrado}) \emph{en un punto}
$x\in X$ si (con las notaciones de (5.9.6)) $\mathcal F_x$ es
$Z_x$-puro (resp. $Z_x$-cerrado); en virtud de (5.9.6), esto equivale a decir
que el homomorfismo canónico
$\mathcal F_x\longrightarrow(\mathcal H^0_{X/Z}(\mathcal F))_x$
es \emph{inyectivo} (resp. \emph{biyectivo}).

Se observará que, para todo $x\in X\setmin Z$, $\mathcal F$ es $Z$-\emph{cerrado en el punto}
$x$, en virtud de (5.9.8).

\begin{corollary}[5.9.10]
\label{IV.2.s5.9.10-fr}
\begin{enumerate}
\item[\textup{(i)}] Sea $(V_\lambda)$ un recubrimiento abierto de $X$. Para que
$\mathcal F$ sea $Z$-puro (\textup{resp.} $Z$-cerrado), es necesario y suficiente que,
para todo $\lambda$, $\mathcal F|V_\lambda$ sea
$(Z\cap V_\lambda)$-puro (\textup{resp.} $(Z\cap V_\lambda)$-cerrado).
\item[\textup{(ii)}] Sea $Z'$ una parte de $Z$ estable por especialización. Si
$\mathcal F$ es $Z$-puro (\textup{resp.} $Z$-cerrado), es $Z'$-puro
(\textup{resp.} $Z'$-cerrado).
\end{enumerate}
\end{corollary}

Esto resulta inmediatamente de (5.9.8, \emph{b')}.
